../cictikz/src/cictikz/data/tex/cictikz_preamble.tex

% The five transistor OTA: NMOS input pair, PMOS current mirror load,
% NMOS tail source. The starting point of every OTA discussion.

\begin{circuitikz}[american, thick, transform shape, circuit ee IEC]

  ../cictikz/src/cictikz/data/tex/cictikz_lib.tex

  \def\xl{0}
  \def\xr{3.2}
  \def\xm{1.6}
  \def\ytail{0}     % tail source at ground
  \def\ysrc{2.2}    % pair source level
  \def\ydr{3.8}     % pair drain level (\ysrc + \grid)
  \def\ytop{5.4}    % supply level (\ydr + \grid)

  % tail transistor
  \draw (\xm,0) \vground;
  \draw (\xm,0) \lvnmos{M5}{$V_{B}$};
  \draw (\xl,\ysrc) -- (\xl,1.9) -- (\xr,1.9) -- (\xr,\ysrc);
  \draw (\xm,\grid) -- (\xm,1.9);
  \fill (\xm,1.9) circle (0.075);

  % input pair
  \draw (\xl,\ysrc) \lvnmos{M1}{$v_p$};
  \draw (\xr,\ysrc) \lvmnmos{M2}{$v_n$};

  % PMOS mirror: M3 diode connected (left), M4 output (right)
  \draw (\xl,\ydr) \lvmpmos{M3}{};
  \draw (\xr,\ydr) \lvpmos{M4}{};
  \draw (\xl,\ytop) \vsupply;
  \draw (\xr,\ytop) \vsupply;

  % gate wire and the diode connection of M3
  \draw (M3.G) -- (M4.G);
  \coordinate (g3) at (M3.G);
  \fill (g3) circle (0.075);
  \draw (g3) |- (\xl,\ydr);
  \fill (\xl,\ydr) circle (0.075);

  % output
  \draw (\xr,\ydr) \portOut{$v_o$};

  % device names
  \node[anchor=west, scale=0.8] at (0.35,3.0) {M1};
  \node[anchor=east, scale=0.8] at (2.85,3.0) {M2};
  \node[anchor=east, scale=0.8] at (-0.35,4.95) {M3};
  \node[anchor=west, scale=0.8] at (3.55,4.95) {M4};
  \node[anchor=west, scale=0.8] at (1.95,0.5) {M5};

\end{circuitikz}
\end{document}
